% The document class supplies options to control rendering of some standard
% features in the result.  The goal is for uniform style, so some attention 
% to detail is *vital* with all fields.  Each field (i.e., text inside the
% curly braces below, so the MEng text inside {MEng} for instance) should 
% take into account the following:
%
% - author name       should be formatted as "FirstName LastName"
%   (not "Initial LastName" for example),
% - supervisor name   should be formatted as "Title FirstName LastName"
%   (where Title is "Dr." or "Prof." for example),
% - degree programme  should be "BSc", "MEng", "MSci", "MSc" or "PhD",
% - dissertation title should be correctly capitalised (plus you can have
%   an optional sub-title if appropriate, or leave this field blank),
% - dissertation type should be formatted as one of the following:
%   * for the MEng degree programme either "enterprise" or "research" to
%     reflect the stream,
%   * for the MSc  degree programme "$X/Y/Z$" for a project deemed to be
%     X%, Y% and Z% of type I, II and III.
% - year              should be formatted as a 4-digit year of submission
%   (so 2014 rather than the academic year, say 2013/14 say).

\documentclass[ % the name of the author
                    author={Daniel Page},
                % the name of the supervisor
                supervisor={Dr. Andrew Calway},
                % the degree programme
                    degree={MSc},
                % the dissertation    title (which cannot be blank)
                     title={Some Structural Guidelines for Data Science MSc Theses, Including Those With Long Titles that Run Across Multiple Lines on the Front Page},
                % the dissertation subtitle (which can    be blank)
                  subtitle={And those including an optional subtitle too, for good measure},
                % the dissertation     type
                      type={},
                % the year of submission
                      year={2021}]{dissertation}

\begin{document}

% =============================================================================

% This section simply introduces the structural guidelines.  It can clearly
% be deleted (or commented out) if you use the file as a template for your
% own dissertation: everything following it is in the correct order to use 
% as is.

\section*{Prelude}
\thispagestyle{empty}


A typical MSc thesis will be structured according to a fairly standard sequence of 
sections, and one example is shown here.  However, it is hard and perhaps even 
counter-productive to generalise: the goal of this document is {\em not}\/ to be prescriptive, to show you how you {\em must} structure your thesis; rather it is to simply to act as a guideline, to show you how you {\em could} structure your thesis.  In particular, each chapter's indicated page-count is
important but {\em not}\/ absolute: the aim is simply to highlight that a 
clear and concise description is better than a rambling alternative that makes
it hard to separate important content and facts from irrelevant trivia.

You can use this document as a \LaTeX-based 
template for your own thesis by simply deleting extraneous sections
and content, and adding your own text. If you do that, we recommend using BibTeX
to deal with the associated bibliography. 
You may want to use Overleaf (an online cloud-based collaborative \LaTeX platform).


The words ``thesis'' and ``dissertation'' are used interchangeably here. All active hyperlinks in the PDF file are rendered in the same shade of dark blue. 



% =============================================================================

% This macro creates the standard UoB title page by using information drawn
% from the document class (meaning it is vital you select the correct degree 
% title and so on).


\maketitle

% After the title page (which is a special case in that it is not numbered)
% comes the front matter or preliminaries; this macro signals the start of
% such content, meaning the pages are numbered with Roman numerals.

\frontmatter

% This macro creates the standard UoB declaration; on the printed hard-copy,
% this must be physically signed by the author in the space indicated.

\makedecl

% LaTeX automatically generates a table of contents, plus associated lists 
% of figures, tables and algorithms.  The former is a compulsory part of the
% dissertation, but if you do not require the latter they can be suppressed
% by simply commenting out the associated macro.

\tableofcontents
\chapter{Context}
\section{Motivations}
1. What chronic pain is 
2. Nociceptors, what they are (pain sensing) and why their activity matters for understanding pain
3. Explain how spike timing and firing patterns of individual fibres carry information about pain
4. Microeletroencephalography, what it is, why it's the right tool, and what it gives. 
5. Identifying individual fibres' spikes in the recording is a problem this thesis addresses 
\section{Challenges}
1. One electrode, many fibres (action potentials), noise
2. Fibre overlap: other fibres fire at similar latencies and have similar waveform shapes to the target. 
3. No clean reference, unlike synthetic data; real recordings have no ground truth clean waveforms, so denoising methods can't be trained or evaluated against a known target. 
\section{Investigative focus}
The aim: a single-channel pipeline to denoise the recording and detect the target unit's spikes and classify them
\section{Structure}
One paragraph that explains what each chapter covers in order
\chapter{Related Work}
\section{Preprocessing and denoising}
1. Standard preprocessing for nerve recordings; bandpass filtering, referencing, artefact blanking 
2. denoising approaches: wavelet methods, autoencoders
3. Challenge specific to spikes, like denoising must preserve spike shape, not just reduce noise
\section{Spike detection and template matching}
1. Amplitude thresholding (the standard k×MAD approach, Quiroga 2004)
2. Template matching (normalised correlation, matching by waveform shape)
3. Latency constraint (using post-stimulus latency to constrain detection).
4. Note the limits of each 
\section{Related pipelines}
1. Troglio (2025 and 2026) hybrid knowledge/data-driven pipeline, latency constraint, SVM classification and their single-electrode 2025 paper.
2. Laboy-Juárez: normalised template matching. but, it's multi-channel (tetrode, template concatenated across 4 electrodes).
3. Buccino / SpikeInterface: quality metrics fail to separate true from false when distributions overlap
4. There's more pipelines 
\section{The gap}

\chapter{Single-Channel Methods}
This chapter describes the data and the methods used to detect, classify and denoise the target unit's spikes in a single channel. It first introduces the recording and how the target unit is identified, then estimates the noise floor. It then describes the detection pipeline in three stages: amplitude thresholding, building and matching a spike template, and constraining detections by their timing after the stimulus. It then describes the two classifiers applied to the surviving detections: a one-class and two-class support vector machine.  Finally, it describes the denoising methods evaluated, the design used to assess whether denoising improves detection or classification, and the metrics used throughout.  

\section{The data: rat saphenous nerve microneurography}
This project uses an extracellular microneurography recording from the peripheral nerve of a rat, sampled at 30 kHz over approximately 2.3 hours. The nerve carries many afferent fibres, so the electrode records the action potentials of many units at once, superimposed on background noise. 
Recordings were made using two multi-electrode sensor units placed at different positions along the nerve, each containing several channels corresponding to slightly different electrode locations. Eight channels carry manually labelled ground truth spikes (CH3, CH8, CH9, CH10, CH14, CH19, CH25, CH30). 
During the recording, an electrical stimulus (Pain stimulus) was delivered repeatedly, giving 2,310 stimulation pulses (TTL) at a normal rate of around 0.25 Hz. The target fibre responds to each stimulus by firing at an approximately fixed latency. The ground truth spike times were marked manually by a human annotator, who only labelled recognisable nerve responses; therefore, the annotations are conservative, and not every response is labelled. Each spike is marked at the point where the signal crosses zero on its descending slope. 
This project develops the pipeline primarily on Channel 19 (CH19), a single channel with a clear, recognisable ground truth. The signal was preprocessed by the recording laboratory, so it has been band-pass filtered 300-6000 Hz, common average referenced, and stimulus blanked. The analysis window used is 900 seconds, starting at the 900th stimulus (~3602.7s), containing 177 ground truth target spikes. 

\section{Noise estimation and preprocessing}
The recording given was already bandpass-filtered, with the stimulation artefact blanked (set to zero); no further filtering was applied in this project. Because the signal is already band-limited, the residual noise is correlated in time rather than white; this will become relevant when evaluating denoising methods later.
The target unit's spikes are small, only about two to three times the amplitude of the background noise, so they do not stand out clearly and can not be isolated by amplitude threshold alone. Therefore, removing noise and detecting these spikes is the central challenge addressed in this chapter. 
The noise level is estimated using the median absolute deviation (MAD) \cite{quiroga_unsupervised_2004, park_deep_2019}:
\begin{equation}
    \sigma_{\text{noise}} = \frac{\operatorname{median}(|x|)}{0.6745}
\end{equation}
where \[x\] is the signal. 
The median is used instead of standard deviation because the spikes are large and infrequent, so the median does not inflate the estimate the way standard deviation would; dividing by 0.6745 rescales the value to be comparable to a standard deviation under an assumption of Gaussian noise \cite{drebitz_optimizing_2019}. MAD is an estimate, so it assumes the signal is mostly noise and approximately Gaussian \cite{coupe_robust_2010}. But because it is applied identically to every signal, it is a valid basis for comparison. This noise estimate forms the basis of the amplitude detection threshold described in the sections below. 

\section{Detection}
\subsection{Amplitude thresholding}
Amplitude thresholding selects spikes whenever the signal drops below a set threshold, so the performance of this method relies mostly on the selected threshold; if the threshold is too low, it will probably recognise many spikes that are not the target; if it is a high threshold, it will lose some of the spikes \cite{cimolato_modern_2021}. (This needs to go to the related work section). 
Spikes are detected from their amplitude. Because the target unit's spikes are negative-going, a detection is registered wherever the signal crosses below a set threshold relative to the noise floor: 
\begin{equation}
    \text{threshold} = -k\,\sigma_{\text{noise}}
    \label{MAD noise estimate}
\end{equation}
Where \(\sigma_{\text{noise}}\) is the median absolute deviation estimate from the previous section and \(k\) is a multiplier controlling how strict the threshold is \cite{quiroga_unsupervised_2004, park_deep_2019}. A downward crossing below this threshold marks a candidate spike. 
To prevent a single spike being counted more than once, a 3 ms refractory period is enforced after each detection, during which no further detection is registered. The reason the refractory period is longer than what is usually reported (0.66 ms in \cite{laboy-juarez_normalized_2019}; 1.28ms in \cite{rizk_optimizing_2009}) because, when threshold crossings were measured without a refractory period, the median interval between consecutive detections was only 1 ms  (95.6\% of intervals fell below 3ms), which is far shorter than the intervals between successive target spikes, which roughly fire once per stimulus (every ~4 s). A single action potential can last several milliseconds and therefore produce several detections milliseconds apart . 


\subsection{Template building and spike detection}
The spike template is built by averaging the ground truth spikes of the target unit. For each marked spike, a 90-sample window (3ms at 30kHz), 45 samples either side of the trough, is then extracted from the signal, and these windows are averaged to give the unit's mean waveform. Before averaging, each spike is aligned on a common reference point within a short search region around the marked time, so that the same feature falls at the same position in every window; without this, the waveform would sit at a slightly different offset, and the average would be blurred
Aligns each spike at its trough (the most negative point), which is a standard and simple approach but is sensitive to noise, since the exact location of the trough can shift slightly from spike to spike \cite{lee_non-linear_2021}. 
To locate spikes matching the target unit's waveform, the template is slid along the signal, and a similarity score is computed at every position. The score is the Pearson correlation coefficient between the template and the signal window at each position, which is an approach by \cite{park_deep_2019} for spike classification on single-electrode extracellular recordings:
\begin{equation}
    r_i = \frac{\sum_{j=1}^{n} (t_j - \bar{t})(x_{i,j} - \bar{x}_i)}
    {\sqrt{\sum_{j=1}^{n} (t_j-\bar{t})^2}\;\sqrt{\sum_{j=1}^{n} (x_{i,j} - \bar{x}_i)^2}}
\end{equation}
where \(t_j\) are the template samples, \(x_i,j\) the signal samples in the window at position \(i,\bar{t}\) and \(\bar{x}\) their mean, and \(n\) is the template length. The score \(r_i\) lies on [-1, 1], where 1 indicates a perfect match in shape, and because the score is normalised, it measures waveform shape independently of amplitude, so a small spike and a large one with the same spike score the same. The principle of using a threshold-normalised shape similarity score to separate the target unit from other detections is used by \cite{laboy-juarez_normalized_2019}, although they apply cosine similarity across a multi-channel (tetrode) template rather than a single-channel Pearson correlation used here. 
 
\subsection{Latency constraint}
Latency-based constraints on the detection search space have been applied in evoked microneurography \cite{troglio_hybrid_2026}. The target fibre fires consistently at about the same latency after each electrical stimulus (~ 100 ms later), so a detection's post-stimulus latency indicated whether it could belong to the target unit or not, which waveform shape alone does not provide. For each ground truth spike, the latency was computed as the interval from the most recent preceding stimulus; the median latency across the 177 target spikes was 100.1 ms. A detection was retained only if its post-stimulus latency fell within an 8 ms window centred on the median of  (96.1-104.1 ms). The window width was selected by sweeping candidates from 1 to 14 ms, and F1 peaked at 8 ms. 
\section{Denoising methods}

 
\subsection{Wavelet}
Wavelet denoising is a signal processing method that works by transforming the signal into wavelet coefficients to remove or shrink the small coefficients that represent noise, and then it reconstructs the signal from what remains \cite{wei_wavelet_2005}. 
For this project, SWT was used to denoise the channel over DWT because SWT does not decimate, so the number of coefficients remains the same across all levels, unlike DWT, which uses a decimation operation; therefore, SWT is better for denoising \cite{kechris_removing_2021}. 
The signal was decomposed using Daubechies 4 to decomposition level 4 to give approximation and details at each level. At each level, the noise level was estimated using the same MAD estimator \ref{MAD noise estimate} on the detail coefficients at each level. A level-independent universal threshold was applied using the VisuShrink formula:
\begin{equation}
    \lambda = \sigma \sqrt{2 \ln N}
    \label{VisuShrink}
\end{equation}
Where \[N\] is the number of coefficients at that level \cite{donoho_minimax_1998}. 
The detail coefficients were shrunk using soft thresholding using the following formula:
\begin{equation}
    \eta(x, \lambda) = \text{sign}(x)\, \max(|x| - \lambda,\, 0)
    \label{Soft thresholding}
\end{equation}
This shrinks the coefficients towards zero instead of removing them. The approximation coefficients were left unchanged. The denoised signal is then reconstructed back from the cleaned coefficients to a time-domain signal by inverse SWT. Soft thresholding was chosen over hard thresholding to allow for smoother reconstruction and to avoid sharp discontinuities\cite{zhang_improved_2019}. 
\subsection{Dense and Convolutional Autoencoders}
Two autoencoder denoisers were used: dense and convolutional autoencoders. Both autoencoders were trained on the same data and applied to the signal in the same way; they only differ in their internal architecture. 
Denoising Autoencoders generally require clean signals as a target reference during training so the model can calculate its reconstructive loss \cite{vincent_extracting_2008}. With real extracellular recordings, there is no clean reference, so training pairs were generated synthetically from the real spike template. The continuous signal was divided into 60 sample windows; each window was created as one of two types: a window that contains the template spike or only a noisy window. For the windows with a spike, the template was randomly shifted and scaled so the model learns to see spikes at different positions and amplitudes. To each clean window, Gaussian noise was then added to produce the noisy input, giving pairs of noisy input and clean target. For training, both autoencoders used the Adam optimiser was to minimise the mean squared error (MSE) between the reconstruction output and the clean target:
\begin{equation}
    \mathcal{L}_{\text{MSE}} = \frac{1}{N} \sum_{i=1}^{N} \left( x_i - \hat{x}_i \right)^2
    \label{MSE}
\end{equation}
Where \[\mathcal{L}_{\text{MSE}}\] is the loss being minimised, \[N\] is the number of samples, \[x_i\] is the clean target value and \[\hat{x}_i\] is the reconstructed output from the Autoencoder \cite{jadon_comprehensive_2022}. 

To denoise the recording, the trained models were applied to overlapping windows sliding across the signal. Each window was normalised by the MAD noise level before the model, to match the training normalisation, then rescaled back, and where windows overlapped, the outputs were averaged.
The difference between Dense and Convolutional autoencoders is that Dense used fully connected layers, compressing the 60-sample window to 30, then to a 15-unit latent space \cite{li_aecuration_2025}. Whereas Convolutional replaces the fully connected layers with 1D convolutional layers, where each convolutional layer slides a small filter (kernel size 3) across the window to detect local spike-shaped patterns wherever they occur in the window \cite{li_accurate_2020, kechris_removing_2021}. This handles spikes at different positions better than dense 

\section{Classification (SVM)}
Two support vector machine (SVM) classifiers with radial basis function (RBF) kernels were applied to trough-aligned waveforms, following the raw waveform approach of \cite{troglio_hybrid_2026}. Both were applied to the latency survivor detections and evaluated by five-fold cross-validation with a fixed random seed. 
[To Do: add either architecture table (comparing the two or something that explains it visually]

\subsection{One-Class SVM}
One-class SVM is a novelty detection machine learning method that learns a decision function to classify new data as similar or different to the training set. Here, the one-class SVM is trained on the target waveforms of latency constraint detection survivors only and evaluated on held-out targets together with all non-target detections. A survivor was labelled a target if it fell within 2 ms of a ground truth spike. The target waveforms were split into five folds; in each fold, the model was trained on the other four folds and tested on the held-out targets together with all non-targets, and it was verified that no waveform appeared in both partitions of any fold. Predictions were pooled across folds so that each candidate received a single out-of-fold prediction, and metrics were computed once on the full set at the true class ratio. An RBF kernel was used (gamma = scale), and nu ($\nu$) was fixed at 0.1 a priori to avoid tuning the hyperparameter on the evaluation data. nu sets the expected fraction of training points allowed to fall outside of the learning boundary. Because the thresholded F1 of a one-class model is sensitive to nu, threshold-independent measures such as ROC-AUC and average precision (PR-AP) were used to summarise performance, as these are suitable for imbalanced classification  \cite{campos_evaluation_2016}. One-class RBF classifiers are sensitive to feature scaling \cite{juszczak_feature_2002, feragen_adaptive_2015}. Standardisation, which rescales each feature by its standard deviation, is problematic for these waveforms because the low-variance tail samples are divided by a very small standard deviation, amplifying them and distorting the distance used by the RBF kernel, so the boundary rejects almost all candidates; therefore, features were not standardised. 
\subsection{Two-class svm}
Unlike one-class SVM, two-class SVM was trained on both target and non-target waveforms to learn the boundary that separates the two classes \cite{worheide_multi-omics_2021}. It was applied to the same latency survivor candidate set as the one-class model. Because targets are a minority of the survivors, class imbalance was handled by class weighting rather than resampling. Features were standardised using a StandardScaler fit on the training folds only and applied to the test fold, in contrast to the one-class model. Similar to one-class, the parameter C was fixed at 1.0 a priori to avoid tuning on the evaluation data.  Evaluation used five-fold stratified cross-validation, with predictions pooled across folds so that each candidate received a single out-of-fold prediction and metrics were computed once on the full set. Similar to one-class, performance was summarised by precision, recall, F1, ROC-AUC and PR-AP. 
[Still I need To DO: There are some fixes I need to do to my code and add bootstrap (single-arm CIs and paired contrast to both one-class and two-class)]


\section{Evaluation metrics}
\section{After denoising evaluation design}

\chapter{Single-Channel Results}
\section{Experimental setup}
\section{Detection performance}
\section{Why detection is limited}
\section{Denoising quality}
\section{The effect of denoising on detection}
\section{Classification performance and the effect of denoising}


\chapter{Discussion}
\section{Critical discussion of findings}
\section{Limitations and recommendations}
\section{Summary}

\chapter{Conclusions and Future Work}
\section{Conclusions}
\section{Recommendations for future work}

\appendix
\chapter{Additional single-channel results}
\chapter{Supplementary methods}
\listoffigures
\listoftables
\listofalgorithms
\lstlistoflistings

% The following sections are part of the front matter, but are not generated
% automatically by LaTeX; the use of \chapter* means they are not numbered.

% -----------------------------------------------------------------------------

\chapter*{Abstract}

{\bf A compulsory section, of at most $1$ page} 
\vspace{1cm} 

\noindent
This section should summarise the project context, aims and objectives,
and main contributions (e.g., deliverables) and achievements.  The goal is to ensure that the 
reader is clear about what the topic is, what you have done within this 
topic, {\em and}\/ {\bf what your view of the outcome is.}

Essentially 
this section is a (very) short version of what is typically covered in more depth in the first 
chapter.  If appropriate, you should include here  
a clear statement of your research hypothesis.  This will obviously differ significantly
for each project, but an example might be as follows:

\begin{quote}
My research hypothesis is that a suitable genetic algorithm will yield
more accurate results (when applied to the standard ACME data set) than 
the algorithm proposed by Jones and Smith, while also executing in less
time.
\end{quote}

\noindent
The latter aspects should (ideally) be presented as a concise, factual 
list of the main points of achievement.  Again the points will differ for each project, but 
an might be as follows:

\begin{quote}
\noindent
\begin{itemize}
\item I spent $120$ hours collecting material on and learning about the 
      Java garbage-collection sub-system. 
\item I wrote a total of $5000$ lines of {\em Python} source code, and associated orchestration scripts. 
\item I designed a new algorithm for computing the non-linear mapping 
      from A-space to B-space using a genetic algorithm.
\item I implemented a version of the algorithm proposed by Jones and 
      Smith (2010), corrected a mistake in it, and 
      compared the results with several alternatives.
\end{itemize}
\end{quote}

% -----------------------------------------------------------------------------

\chapter*{Summary of Changes}

{\bf A conditional section, of at most $1$ page} 
\vspace{1cm} 

If (and only if) the dissertation represents a resubmission (e.g., as the result of
a resit), then this section is compulsory: the content should summarise all
non-trivial changes made to the initial submission.  Otherwise you can
omit it, since {\bf a summary of this type is only needed for resubmissions}.

When included, the section will ideally be used to highlight additional
work completed, and address criticism raised in any associated feedback.
Clearly it is difficult to give generic advice about how to do so, but
an example might be as follows:

\begin{quote}
\noindent
\begin{itemize}
\item Feedback from the initial submission criticised the design and 
      implementation of my genetic algorithm, stating ``there seems 
      to have been no attention to computational complexity during the
      design, and obvious methods of optimisation are missing within
      the resulting implementation''.  Chapter 3 now includes a
      comprehensive analysis of the algorithm, in terms of both time
      and space.  While I have not altered the algorithm itself, I
      have included a cache mechanism (also detailed in Chapter 3)
      that provides a significant improvement in average run-time.
\item I added a feature in my implementation to allow automatic rather
      than manual selection of various parameters; the experimental
      results in Chapter 4 have been updated to reflect this.
\item Questions after the presentation highlighted a range of related
      work that I had not considered: I have make a number of updates 
      to Chapter 2, resolving this issue.
\end{itemize}
\end{quote}

% -----------------------------------------------------------------------------

\chapter*{Supporting Technologies}

{\bf A compulsory section, of at most $1$ page}
\vspace{1cm} 

\noindent
This section should present a detailed summary, in bullet point form, 
of any third-party resources (e.g., hardware and software components) 
used during the project.  Use of such resources is always perfectly 
acceptable: the goal of this section is simply to be clear about how
and where they are used, so that a clear assessment of your work can
result.  The content can focus on the project topic itself (rather,
for example, than including ``I used \mbox{\LaTeX} to prepare my 
dissertation''); an example is as follows:

\begin{quote}
\noindent
\begin{itemize}

\item I used the {\em Pandas} and {\em Seaborn} public-domian Python Libraries. 

\item I used a parts of the OpenCV computer vision library to capture 
      images from a camera, and for various standard operations (e.g., 
      threshold, edge detection).

\item I used Amazon Web Services for remote storage and processing of data. Specifically, I used:
 \begin{itemize}
 \item Simple Storage Service (S3) for data storage
 \item Elastic Compute Cloud (EC2) for provision of virtual machines
 \item Elastic Beanstalk for scaling and load management
 \item Sagemaker for all the machine learning components of my project. 
 \end{itemize}
 
\item I used \LaTeX\ to format my thesis, via the online service {\em Overleaf}. 
\end{itemize}
\end{quote}

% -----------------------------------------------------------------------------

\chapter*{Notation and Acronyms}

{\bf An optional section, of roughly $1$ or $2$ pages}
\vspace{1cm} 

\noindent
Any well written document will introduce notation and acronyms before
their use, {\em even if} they are standard in some way: this ensures 
any reader can understand the resulting self-contained content.  

Said introduction can exist within the dissertation itself, wherever 
that is appropriate.  For an acronym, this is typically achieved at 
the first point of use via ``Advanced Encryption Standard (AES)'' or 
similar, noting the capitalisation of relevant letters.  However, it 
can be useful to include an additional, dedicated list at the start 
of the dissertation; the advantage of doing so is that you cannot 
mistakenly use an acronym before defining it.  A limited example is 
as follows:

\begin{quote}
\noindent
\begin{tabular}{lcl}
AES                 &:     & Advanced Encryption Standard                                         \\
DES                 &:     & Data Encryption Standard                                             \\
                    &\vdots&                                                                      \\
${\mathcal H}( x )$ &:     & the Hamming weight of $x$                                            \\
${\mathbb  F}_q$    &:     & a finite field with $q$ elements                                     \\
$x_i$               &:     & the $i$-th bit of some binary sequence $x$, st. $x_i \in \{ 0, 1 \}$ \\
\end{tabular}
\end{quote}

% -----------------------------------------------------------------------------

\chapter*{Acknowledgements}

{\bf An optional section, of at most $1$ page}
\vspace{1cm} 

\noindent
It is common practice (although totally optional) to acknowledge any
third-party advice, contribution or influence you have found useful
during your work.  Examples include support from friends or family, 
the input of your Supervisor and/or Advisor, external organisations 
or persons who  have supplied resources of some kind (e.g., funding, 
advice or time), and so on.

\vspace{1cm}
Dave Cliff writes here to say huge thanks to his colleague Dr Dan Page for sharing this \LaTeX\ thesis template, which was originally written by Dan, for Computer Science dissertations. Dave edited Dan's original to better suit the needs of the Data Science MSc: please don't hassle Dan about any of this, but do feel free to contact Dave if you have any questions or comments on it.  

% =============================================================================

% After the front matter comes a number of chapters; under each chapter,
% sections, subsections and even subsubsections are permissible.  The
% pages in this part are numbered with Arabic numerals.  Note that:
%
% - A reference point can be marked using \label{XXX}, and then later
%   referred to via \ref{XXX}; for example Chapter\ref{chap:context}.
% - The chapters are presented here in one file; this can become hard
%   to manage.  An alternative is to save the content in seprate files
%   the use \input{XXX} to import it, which acts like the #include
%   directive in C.

\mainmatter

% -----------------------------------------------------------------------------

\chapter{Introduction}
\label{chap:introduction}

{\bf A compulsory chapter, roughly 10\% of the total page-count}
\vspace{1cm} 

% putting a \noindent before the first para in each chapter looks nicer.
\noindent
This chapter should describe the project context, and motivate each of
the proposed aims and objectives.  Ideally, it is written at a fairly 
high-level, and easily understood by a reader who is technically 
competent but not an expert in the topic itself.

In short, the goal is to answer three questions for the reader.  First, 
what is the project topic, or problem being investigated?  Second, why 
is the topic important, or rather why should the reader care about it?  
For example, why there is a need for this project, who will benefit from the 
project and in what way (e.g., clients/end-users who needed some analysis
done, or other data scientists who might need the tools you have developed), what 
work does the project build on and why is the selected approach either
important and/or interesting (e.g., fills a gap in literature, applies
results from another field to a new problem).  Finally, what are the 
central challenges involved and why are they significant? 
 
The chapter should conclude with a concise bullet point list that 
summarises the aims, objectives, {\bf and achievements}\/ of your work. 

% -----------------------------------------------------------------------------

\chapter{Background}
\label{chap:background}

{\bf A compulsory chapter, roughly 20\% of the total page-count}
\vspace{1cm} 

\noindent
This chapter is intended to describe the technical basis on which execution
of the project depends.  The goal is to provide a detailed explanation of
the specific problem at hand, and existing work that is relevant such as an
existing algorithm that you use, alternative solutions proposed, supporting
technologies, and relevant literature. The literature you include should cover 
appropriate peer-reviewed academic publications and books, and maybe also 
news and current-affairs articles published in reputable sources such as 
{\em The Economist} magazine or {\em The Financial Times} newspaper. 

Your thesis should include complete set of references/bibliography: everything that you cite should be in there, in full. This means including the publisher's name for anything that is a book; the editors and title of books that a paper appears in as one item in a larger collection (e.g. proceedings volumes and/or thematic edited collections) and the relevant page-numbers; and so on. Put most simply, we require and expect your references to look like the references in a published peer-reviewed academic paper, so you can work out whether you still have work to do by comparing your references to the reference-list on the publications that you're working from. 

Note there is a subtle difference from
this and a full-blown literature {\em survey}.  The latter might try
to capture and organise (e.g., categorise somehow) {\em all}\/ related work,
potentially offering meta-analysis, whereas here the goal is simply to
ensure that your thesis is self-contained.  Put another way, after reading 
this chapter an intelligent non-expert reader with no prior knowledge of your project should have obtained enough background to 
understand what {\em you}\/ have done (by reading subsequent sections), and then 
accurately assess your work.  You might view an additional goal as giving 
the reader confidence that you are able to absorb, understand and clearly 
communicate highly technical material.

Just as there is no single ideal structure for an MSc thesis, there is no one correct name for this chapter. You could just call it {\em Background}\/ if you wish, or {\em Technical Background}. Or if you feel you have a lot to say you could split this chapter into two: you might have your first three chapters being:
\begin{enumerate}
\item Introduction
\item Context
\item Related Work
\end{enumerate}
\noindent
But if you prefer to use a more concise structure, you could instead use:
\begin{enumerate}
\item Introduction: Contextual Background
\item Technical Background
\end{enumerate}
\noindent
The choice is yours. 


The key thing is that by the end of these first two or three chapters you have told the reader everything they need to know so that they can understand the rest of your thesis. This is particularly important because at least one of the people who actually examines your thesis will be a UoB academic, one of our lecturers or professors, who you can assume knows almost nothing about the specifics of what work you have done, and who is given a copy of your thesis to mark. So what you write has to adequately explain things to that examiner. You can safely assume that whoever examines your thesis is numerate, intelligent, understands programming and data analytics etc, but you cannot safely assume that their specific individual expertise is a perfect match to the topic of your thesis (it's in that sense that the examiner is a {\em non-expert}). So, the person you write for, your {\em audience}, should be that undefined group of academics who might possibly examine your thesis: if you write for that audience and do a good job, your thesis should be understandable by a wide range of people, including potential employers and colleagues in the world of work.

% -----------------------------------------------------------------------------

\chapter{Execution}
\label{chap:execution}

{\bf A topic-specific chapter, roughly 30\% of the total page-count}
\vspace{1cm} 

\noindent
This chapter is intended to describe what you did: the goal is to explain
the main activity or activities, of any type, which constituted your work 
during the project.  The content is highly topic-specific. For some 
projects it will make sense to split the content into two main sections, or maybe even into two separate chapters: one 
will discuss the design of something, including any rationale or decisions made, 
and the other will discuss how this design was realised via some form of 
implementation.  You could instead give this chapter the title ``Design and Implementation"; or you might split this content into two chapters, one titled ``Design" and the other ``Implementation" .

Note that it is common to include evidence of ``best practice'' project 
management (e.g., use of version control, choice of programming language 
and so on).  Rather than simply a rote list, make sure any such content 
is useful and/or informative in some way: for example, if there was a 
decision to be made then explain the trade-offs and implications 
involved.

\section{Example Section}

This is an example section; 
the following content is auto-generated dummy text.
\lipsum

\subsection{Example Sub-section}

\begin{figure}[t]
\centering
foo
\caption{This is an example figure.}
\label{fig}
\end{figure}

\begin{table}[t]
\centering
\begin{tabular}{|cc|c|}
\hline
foo      & bar      & baz      \\
\hline
$0     $ & $0     $ & $0     $ \\
$1     $ & $1     $ & $1     $ \\
$\vdots$ & $\vdots$ & $\vdots$ \\
$9     $ & $9     $ & $9     $ \\
\hline
\end{tabular}
\caption{This is an example table.}
\label{tab}
\end{table}

\begin{algorithm}[t]
\For{$i=0$ {\bf upto} $n$}{
  $t_i \leftarrow 0$\;
}
\caption{This is an example algorithm.}
\label{alg}
\end{algorithm}

\begin{lstlisting}[float={t},caption={This is an example listing.},label={lst},language=C]
for( i = 0; i < n; i++ ) {
  t[ i ] = 0;
}
\end{lstlisting}

This is an example sub-section;
the following content is auto-generated dummy text.
Notice the examples in Figure~\ref{fig}, Table~\ref{tab}, Algorithm~\ref{alg}
and Listing~\ref{lst}.
\lipsum

\subsubsection{Example Sub-sub-section}

This is an example sub-sub-section;
the following content is auto-generated dummy text.
\lipsum

\paragraph{Example paragraph.}

This is an example paragraph; note the trailing full-stop in the title,
which is intended to ensure it does not run into the text.

% -----------------------------------------------------------------------------

\chapter{Critical Evaluation}
\label{chap:evaluation}

{\bf A topic-specific chapter, roughly 30\% of the total page-count} 
\vspace{1cm} 

\noindent
This chapter is intended to evaluate what you did.  The content is highly 
topic-specific, but for many projects will have flavours of the following:

\begin{enumerate}
\item functional  testing, including analysis and explanation of failure 
      cases,
\item behavioural testing, often including analysis of any results that 
      draw some form of conclusion wrt. the aims and objectives,
      and
\item evaluation of options and decisions within the project, and/or a
      comparison with alternatives.
\end{enumerate}

\noindent
This chapter often acts to differentiate project quality: even if the work
completed is of a high technical quality, critical yet objective evaluation 
and comparison of the outcomes is crucial.  In essence, the reader wants to
learn something, so the worst examples amount to simple statements of fact 
(e.g., ``graph X shows the result is Y''); the best examples are analytical 
and exploratory (e.g., ``graph X shows the result is Y, which means Z; this 
contradicts [1], which may be because I use a different assumption'').  As 
such, both positive {\em and}\/ negative outcomes are valid {\em if} presented 
in a suitable manner.

% -----------------------------------------------------------------------------

\chapter{Conclusion}
\label{chap:conclusion}

{\bf A compulsory chapter,  roughly 10\% of the total page-count}
\vspace{1cm} 

\noindent
The concluding chapter(s) of a dissertation are often underutilized because they're 
too often left too close to the deadline: it is important to allocate enough time and 
attention to closing off the story, the narrative, of your thesis.

Again, there is no single correct way of closing a thesis. 

One good way of doing this is to have a single chapter consisting of three parts:

\begin{enumerate}
\item (Re)summarise the main contributions and achievements, in essence
      summing up the content.
\item Clearly state the current project status (e.g., ``X is working, Y 
      is not'') and evaluate what has been achieved with respect to the 
      initial aims and objectives (e.g., ``I completed aim X outlined 
      previously, the evidence for this is within Chapter Y'').  There 
      is no problem including aims which were not completed, but it is 
      important to evaluate and/or justify why this is the case.
\item Outline any open problems or future plans.  Rather than treat this
      only as an exercise in what you {\em could} have done given more 
      time, try to focus on any unexplored options or interesting outcomes
      (e.g., ``my experiment for X gave counter-intuitive results, this 
      could be because Y and would form an interesting area for further 
      study'' or ``users found feature Z of my software difficult to use,
      which is obvious in hindsight but not during at design stage; to 
      resolve this, I could clearly apply the technique of Bloggs {\em et al.}.
\end{enumerate}

Alternatively, you might want to divide this content into two chapters: a penultimate chapter with a title such as ``Further Work" and then a final chapter ``Conclusions". Again, there is no hard and fast rule, we trust you to make the right decision. 

And this, the final paragraph of this thesis template, is just a bunch of citations, added to show how to generate a BibTeX bibliography. Sources that have been randomly chosen to be cited here include:
\cite{miller_etal_2018_clojure,webber_marwan_2015,touretzky_2013_lisp,eckmann_etal_1987,marwan_2011,vach_2015,shiller_2017,vytelingum_2006,tesfatsion_2002,rust_etal_1992}.




% =============================================================================

% Finally, after the main matter, the back matter is specified.  This is
% typically populated with just the bibliography.  LaTeX deals with these
% in one of two ways, namely
%
% - inline, which roughly means the author specifies entries using the 
%   \bibitem macro and typesets them manually, or
% - using BiBTeX, which means entries are contained in a separate file
%   (which is essentially a database) then imported; this is the 
%   approach used below, with the databased being dissertation.bib.
%
% Either way, the each entry has a key (or identifier) which can be used
% in the main matter to cite it, e.g., \cite{X}, \cite[Chapter 2}{Y}.

\backmatter

\bibliography{references}

% -----------------------------------------------------------------------------

% The dissertation concludes with a set of (optional) appendices; these are 
% the same as chapters in a sense, but once signalled as being appendices via
% the associated macro, LaTeX manages them appropriately.

\appendix

\chapter{An Example Appendix}
\label{appx:example}

Content which is not central to, but may enhance the dissertation can be 
included in one or more appendices; examples include, but are not limited
to

\begin{itemize}
\item lengthy mathematical proofs, numerical or graphical results which 
      are summarised in the main body,
\item sample or example calculations, 
      and
\item results of user studies or questionnaires.
\end{itemize}

\noindent
Note that in line with most research conferences, the examiners are not
obliged to read such appendices.

% =============================================================================

\end{document}
